% =====================================================================
% 第 5 章  实证结果
% ---------------------------------------------------------------------
% 编译说明：
%   主文档建议用 ctexart（\documentclass[12pt]{ctexart}），本文件作为
%   % =====================================================================
% 第 5 章  实证结果
% ---------------------------------------------------------------------
% 编译说明：
%   主文档建议用 ctexart（\documentclass[12pt]{ctexart}），本文件作为
%   % =====================================================================
% 第 5 章  实证结果
% ---------------------------------------------------------------------
% 编译说明：
%   主文档建议用 ctexart（\documentclass[12pt]{ctexart}），本文件作为
%   % =====================================================================
% 第 5 章  实证结果
% ---------------------------------------------------------------------
% 编译说明：
%   主文档建议用 ctexart（\documentclass[12pt]{ctexart}），本文件作为
%   \input{ch5_empirical} 引入；若主文档是 ctexbook/report，把下面的
%   \section{实证结果} 改成 \chapter{实证结果}，并把 \subsection 降一级。
%   依赖宏包：booktabs、threeparttable、amsmath、graphicx
%   表格编号按"第 5 章"手工写定，若章节顺序调整需同步修改。
% =====================================================================

\section{实证结果}

本章报告文本情绪对市场收益与波动的影响。基准样本为 2001 年第一季度至
2026 年第二季度共 102 期《货币政策执行报告》，其中 99 期可完整切分出政策指引段
与宏观回顾段的情绪指标。市场收益由 16 个连续可得的申万一级行业等权组合构造，
并以 28 个行业口径作为稳健性对照。情绪变量均经标准化处理，回归系数因此可解释为
情绪每变动一个标准差的效应；标准误采用 Newey--West 异方差自相关稳健估计，
滞后阶数取 4。

\subsection{描述性统计}

表~\ref{tab:desc} 列出主要变量的分布。政策指引段情绪（tone\_P）与宏观回顾段情绪
（tone\_R）的相关系数为 0.375，二者相关但并非同一信息的重复度量，这为后续的
分段检验提供了必要的前提。事件日市场收益的均值为 19 个基点，标准差 186 个基点；
日内振幅均值为 2.09\%，意味着报告发布当日最高价与最低价之差平均约为前收盘价的
2\%。事前波动（发布前 60 个交易日市场收益的标准差）均值为 1.46\%。

\begin{table}[htbp]
\centering
\small
\begin{threeparttable}
\caption{描述性统计}
\label{tab:desc}
\begin{tabular}{lcccccc}
\toprule
变量 & 均值 & 标准差 & 最小值 & 中位数 & 最大值 & 观测数 \\
\midrule
政策指引段情绪（tone\_P，标准化） & 0.000 & 1.005 & $-2.773$ & 0.029 & 2.558 & 99 \\
宏观回顾段情绪（tone\_R，标准化） & 0.000 & 1.005 & $-2.928$ & 0.068 & 2.535 & 102 \\
发布后一日收益（\%）              & 0.19  & 1.86  & $-7.60$ & 0.30  & 4.63  & 102 \\
发布后五日收益（\%）              & 0.52  & 3.59  & $-12.59$ & 0.71 & 11.58 & 102 \\
事件日日内振幅（\%）              & 2.09  & 0.92  & 0.80  & 1.88  & 4.76  & 101 \\
事件日绝对收益（\%）              & 0.97  & 0.92  & 0.00  & 0.66  & 3.70  & 102 \\
事前 60 日波动（\%）              & 1.46  & 0.64  & 0.61  & 1.33  & 3.51  & 102 \\
\bottomrule
\end{tabular}
\begin{tablenotes}[flushleft]
\footnotesize
\item 样本为 2001Q1--2026Q2 共 102 期报告；收益与振幅由 16 个申万一级行业的等权组合计算。
\end{tablenotes}
\end{threeparttable}
\end{table}

值得先说明的是情绪与市场波动状态的原始相关性。tone\_P 与事件日日内振幅的相关系数为
$-0.302$，与事件日绝对收益为 $-0.254$，与事前波动的相关系数为 $-0.232$。最后一个
数字提示，情绪较高的报告往往出现在市场相对平静的时期，若不加以控制，很容易把
``情绪降低波动''与``情绪在高波动期偏低''混为一谈。因此在回归中同时报告含事前波动
控制的设定，并以此作为主要判断依据。

\subsection{政策指引情绪与市场波动}

表~\ref{tab:vol} 是本章的主结果。第 (1) 列以事件日绝对收益衡量波动，tone\_P 的系数为
$-23.4$ 个基点（$t=-2.53$）；加入事前波动后（第 (2) 列），系数降至 $-16.6$，显著性
减弱至 10\% 水平，说明有一部分关联确实来自波动状态的持续性。由于绝对收益无法区分
``价格单边移动''与``日内大幅震荡''，第 (3)、(4) 列改用日内振幅，系数分别为 $-27.8$
（$t=-2.58$）与 $-18.5$（$t=-2.43$），控制事前波动后依然在 5\% 水平显著。

第 (5) 列同时纳入回顾段情绪。tone\_R 的系数为 $-2.1$（$t=-0.20$），几乎为零，而
tone\_P 维持在 $-17.8$（$t=-2.27$）。这一对比有明确含义：显著影响波动的并不是
``整篇报告听起来是否乐观''，而是其中面向未来的政策指引部分。第 (6) 列把振幅扩展到
发布后五个交易日，系数为 $-12.1$（$t=-2.03$），效应并未随窗口延长而累积，更接近于
一次性的信息释放而非持续影响。

就经济意义而言，以振幅均值 2.09\% 为参照，$-18.5$ 个基点相当于振幅下降约 9\%；
以标准差衡量，一个标准差的情绪变动对应约 0.2 个标准差的振幅变化。这一量级不算大，
但在 99 个事件、日频噪声显著的条件下已属不易识别。

需要提醒的是，这一幅度对是否包含危机年份较为敏感。剔除 2007--2010 年后，系数降至
$-15.5$（$t=-2.07$），约为基准的一半（见 5.5 节）。方向与显著性不变，但把 $-28$
读作一个稳定不变的效应量并不恰当，更稳妥的表述是：效应量在 15 至 28 个基点之间，
取决于样本是否包含金融危机时期。

分组视角补充了一个细节。把情绪按三分位分组后，低情绪组的平均振幅为 2.53\%，中、高
两组分别为 1.75\% 与 1.99\%（图 5.2）。若把样本切成五组，均值依次为
2.83\%、2.07\%、1.65\%、2.13\%、1.78\%，同样不是单调下降。也就是说，效应主要来自
情绪偏低的一端：措辞异常谨慎的时期波动显著更高，而中高情绪之间并无梯度。这与
``沟通收敛分歧''的解释相容——市场在意的未必是措辞有多乐观，而是它是否透露出不安。

\begin{table}[htbp]
\centering
\small
\begin{threeparttable}
\caption{政策指引段情绪与市场波动}
\label{tab:vol}
\begin{tabular}{lcccccc}
\toprule
 & (1) & (2) & (3) & (4) & (5) & (6) \\
 & $|r|$ & $|r|$ & 振幅 & 振幅 & 振幅 & 五日振幅 \\
\midrule
tone\_P        & $-23.4^{**}$ & $-16.6^{*}$ & $-27.8^{***}$ & $-18.5^{**}$ & $-17.8^{**}$ & $-12.1^{**}$ \\
                & $(-2.53)$ & $(-1.96)$ & $(-2.58)$ & $(-2.43)$ & $(-2.27)$ & $(-2.03)$ \\
tone\_R        & & & & & $-2.1$ & \\
                & & & & & $(-0.20)$ & \\
事前波动        & & 46.0$^{***}$ & & 61.9$^{***}$ & 61.1$^{***}$ & 80.0$^{***}$ \\
                & & $(4.10)$ & & $(4.24)$ & $(4.14)$ & $(4.16)$ \\
\midrule
观测数          & 99 & 99 & 98 & 98 & 98 & 99 \\
$R^2$           & 0.065 & 0.162 & 0.091 & 0.267 & 0.268 & 0.355 \\
\bottomrule
\end{tabular}
\begin{tablenotes}[flushleft]
\footnotesize
\item 被解释变量为事件日市场波动，单位为百分比；情绪与事前波动的系数均已换算为基点
（1 个基点 $=0.01$ 个百分点），其中事前波动的系数读作``事前波动每上升 1 个百分点''。
括号内为 Newey--West 稳健 $t$ 值。$^{*}$、$^{**}$、$^{***}$ 分别表示 10\%、5\%、1\% 显著性水平。
\end{tablenotes}
\end{threeparttable}
\end{table}

\subsection{制度断点：2014 年前后}

样本期内货币政策框架发生了明显变化：2013 年流动性紧张、2014 年利率市场化加速、
2015 年汇率形成机制调整，报告自身的章节结构也在 2015 年前后改写。这些变化提示
情绪的作用方式未必稳定。表~\ref{tab:break} 在方程中引入 2014 年虚拟变量及其与
情绪的交互项，直接检验这一可能。

收益方程中，交互项系数为 $-75.6$ 个基点（$t=-2.11$），五日窗口扩大到 $-181.4$
（$t=-2.65$），即 2014 年以前情绪对收益方向无显著影响（点估计 $+31.6$，$t=0.97$），
2014 年以后转为负向。波动方程则呈现相反的变化：2014 年以前的系数为 $-35.0$
（$t=-3.10$），交互项为 $+26.6$（$t=1.59$），说明波动渠道在后期减弱，但未完全消失。

需要说明的是断点的选择并非由数据挑选而来，而是对应上述制度节点。为检验这一处理是否
稳健，表~\ref{tab:breaksens} 将断点依次替换为 2008、2010、2012、2014、2016、2018 年。
2008 年作断点时交互项接近于零，其余各年份的系数方向一致，其中 2010、2012、2014、2018
年显著，2016 年偏弱。这说明变化是渐进的，把它归因于某一个特定年份并不合适，但
``2010 年代之后效应改变''这一判断是站得住的。

\begin{table}[htbp]
\centering
\small
\begin{threeparttable}
\caption{制度断点检验}
\label{tab:break}
\begin{tabular}{lcccc}
\toprule
 & (1) & (2) & (3) & (4) \\
 & 一日收益 & 五日收益 & 日内振幅 & $|r|$ \\
\midrule
tone\_P              & $+31.6$   & $+73.7$   & $-35.0^{***}$ & $-32.7^{***}$ \\
                      & $(0.97)$  & $(1.22)$  & $(-3.10)$     & $(-2.97)$ \\
tone\_P $\times$ 2014后 & $-75.6^{**}$ & $-181.4^{***}$ & $+26.6$ & $+25.1$ \\
                      & $(-2.11)$ & $(-2.65)$ & $(1.59)$      & $(1.53)$ \\
\midrule
观测数                & 99 & 99 & 98 & 99 \\
$R^2$                 & 0.076 & 0.077 & 0.190 & 0.084 \\
\bottomrule
\end{tabular}
\begin{tablenotes}[flushleft]
\footnotesize
\item 各列均控制 2014 年虚拟变量本身；系数单位为基点，括号内为 Newey--West 稳健 $t$ 值。
\end{tablenotes}
\end{threeparttable}
\end{table}

\begin{table}[htbp]
\centering
\small
\begin{threeparttable}
\caption{断点年份敏感性（交互项系数与 $t$ 值）}
\label{tab:breaksens}
\begin{tabular}{lccc}
\toprule
断点年份 & 一日收益 & 五日收益 & 日内振幅 \\
\midrule
2008 & $-5.9\ (-0.14)$   & $-1.6\ (-0.02)$    & $+1.8\ (0.10)$ \\
2010 & $-86.6^{**}\ (-2.41)$ & $-193.1^{***}\ (-2.69)$ & $+34.5^{**}\ (2.04)$ \\
2012 & $-76.6^{**}\ (-2.20)$ & $-183.5^{***}\ (-2.74)$ & $+27.4^{*}\ (1.71)$ \\
2014 & $-75.6^{**}\ (-2.11)$ & $-181.4^{***}\ (-2.65)$ & $+26.6\ (1.59)$ \\
2016 & $-57.7\ (-1.47)$  & $-134.9^{*}\ (-1.88)$   & $+34.4^{**}\ (2.07)$ \\
2018 & $-64.4^{*}\ (-1.82)$ & $-164.9^{**}\ (-2.29)$ & $+16.2\ (0.79)$ \\
\bottomrule
\end{tabular}
\begin{tablenotes}[flushleft]
\footnotesize
\item 每行对应一次单独的回归，仅改变断点年份，其余设定与表~\ref{tab:break} 相同。
\end{tablenotes}
\end{threeparttable}
\end{table}

\subsection{2014 年之后的方向性效应}

既然断点检验显示后期情绪与收益转为负相关，这里单独报告 2014 年以后的子样本。
表~\ref{tab:post14} 显示，tone\_P 对发布后一日收益的系数为 $-44.0$ 个基点
（$t=-2.89$），五日窗口为 $-107.8$（$t=-3.30$）；同时，波动类方程在子样本中
全部不显著。两组结果互为印证：后期的定价体现在方向上，而不在波动上。

对这一结果需要保持克制。首先，负号的含义并非``情绪越宽松市场越差''这么简单，
更可能反映宽松措辞集中出现在基本面走弱的时点，市场将政策宽松解读为经济前景的
负面信号。其次，该结果对危机年份高度敏感：剔除 2015 年后，一日窗口的交互项由
$-75.6$（$t=-2.11$）降至 $-59.0$（$t=-1.62$）；剔除 2008--2009 年后降至 $-45.5$
（$t=-1.56$）。任一段被剔除，5\% 水平的显著性即不复存在。因此本文不把它作为独立
结论，仅记录符号的方向性变化。

\begin{table}[htbp]
\centering
\small
\begin{threeparttable}
\caption{2014 年以后的子样本}
\label{tab:post14}
\begin{tabular}{lcccc}
\toprule
 & (1) & (2) & (3) & (4) \\
 & 一日收益 & 五日收益 & 一日·双段 & 日内振幅 \\
\midrule
tone\_P & $-44.0^{***}$ & $-107.8^{***}$ & $-47.2^{**}$ & $-8.4$ \\
        & $(-2.89)$     & $(-3.30)$      & $(-2.36)$    & $(-0.68)$ \\
tone\_R &               &                & $+5.9$       & \\
        &               &                & $(0.28)$     & \\
\midrule
观测数  & 50 & 50 & 50 & 49 \\
\bottomrule
\end{tabular}
\begin{tablenotes}[flushleft]
\footnotesize
\item 子样本为 2014Q1--2026Q2。系数单位为基点，括号内为 Newey--West 稳健 $t$ 值。
\end{tablenotes}
\end{threeparttable}
\end{table}

\subsection{稳健性检验}

表~\ref{tab:robust} 汇总四个方面的检验，表~\ref{tab:robust2} 补充三项针对性检验。

\emph{市场口径}。把等权组合由 16 个行业扩展为 28 个行业，主结果基本不变：
日内振幅系数 $-26.7$（$t=-2.58$），断点交互项 $-73.1$（$t=-2.01$），五日交互项
$-183.6$（$t=-2.71$）。结论不依赖于行业集合的选择。

\emph{词典}。以本文自建的货币政策领域词典（积极 159 词、消极 139 词）重新度量情绪，
各系数方向与主结果一致，但量级缩小、显著性下降：绝对收益方程仍显著（$-19.1$，
$t=-2.43$），日内振幅方程降至 $-7.8$（$t=-0.66$），断点交互项降至 $-60.5$
（$t=-1.25$）。该词典每千字仅命中约 68 个情感词，主词典为 104 个，覆盖率差异恰好
对应系数的等比例收缩，符合测量误差导致衰减的预期。因此这一结果是``方向稳健、精度不足''，
而非``结论被推翻''。

\emph{影响点}。逐一剔除 2015 年（股市异常波动）与 2020 年（疫情冲击），并对收益做
1\%/99\% 缩尾。波动类结果对上述处理均不敏感，日内振幅系数稳定在 $-25.8$ 至 $-27.9$
之间。断点交互项则较为敏感：剔除 2015 年后，一日窗口的交互项由 $-75.6$（$t=-2.11$）
降至 $-59.0$（$t=-1.62$），显著性消失；五日窗口降至 $-152.5$（$t=-2.29$），仍在 5\%
水平显著。剔除 2020 年与收益缩尾后，一日交互项分别为 $-73.0$ 与 $-74.4$，保持在 5\%
水平。可见后期方向效应确有相当一部分由 2015 年的数期报告贡献，这解释了上文为何
采取保留态度，也说明该效应在更长窗口上比在单日窗口上更稳定。

\emph{置换检验}。将情绪序列与事件随机配对 500 次，观察真实 $t$ 值在随机分布中的位置。
日内振幅方程的置换 $p$ 值为 0.020，加入回顾段与事前波动后为 0.068，断点交互项为 0.058。
置换检验不依赖分布假设，其结果与稳健 $t$ 检验基本一致：波动结果在 5\% 水平成立，
断点结果处于边际。

\emph{危机期敏感性}。2007--2009 年的措辞与波动同时处于极端，是本文最担心的混淆来源。
剔除 2007--2010 年后，振幅方程系数降至 $-15.5$（$t=-2.07$），剔除 2008--2009 年后为
$-18.8$（$t=-2.34$），加入年份趋势后为 $-20.7$（$t=-2.18$）。方向与显著性保持不变，
幅度约为基准的一半。断点交互项则不同：剔除 2008--2009 年后降至 $-45.5$（$t=-1.56$），
不再显著。可见波动结果稳健但幅度受危机年份影响，方向结果则更依赖少数极端时期。

\emph{安慰剂窗口}。把事件日整体平移 $\pm$5、$\pm$10、$\pm$20 个交易日重估。平移 $\pm$10
天以上时，振幅方程的系数不再显著（$t$ 值均小于 1.1），断点交互项也不再成立，说明效应
与发布时点相关。但平移 $-5$ 天时系数仍有 $-26.6$（$t=-1.79$），保留了约九成的幅度。
这表明效应并非严格的单日冲击，而是集中在发布前后约一周的窗口内，与报告部分内容提前
释放的解释相容。

\emph{政策操作控制}。一个自然的担心是，文本情绪只是货币政策操作的代理。本文整理了三类
总量操作——全面降准 24 次、1 年期 LPR 下调 12 次、7 天逆回购利率下调 12 次——按报告期
汇总为季度操作强度，加入振幅方程。情绪系数由 $-27.8$ 变为 $-27.3$（$t=-2.53$），加入
事前波动后为 $-16.7$（$t=-2.17$）；三类操作本身均不显著，与情绪的相关性也只在 $-0.01$
到 $-0.21$ 之间。断点交互项由 $-75.5$ 变为 $-73.9$（$t=-2.12$）。这说明波动效应并非
由降准降息驱动。

\begin{table}[htbp]
\centering
\small
\begin{threeparttable}
\caption{稳健性检验汇总}
\label{tab:robust}
\begin{tabular}{lcccc}
\toprule
 & 主结果 & 28 行业 & 替代词典 & 剔除 2015 \\
\midrule
\multicolumn{5}{l}{\emph{面板 A：日内振幅（系数，$t$ 值）}} \\
tone\_P & $-27.8^{***}$ & $-26.7^{***}$ & $-7.8$ & $-25.8^{**}$ \\
        & $(-2.58)$ & $(-2.58)$ & $(-0.66)$ & $(-2.19)$ \\
\midrule
\multicolumn{5}{l}{\emph{面板 B：断点交互项（系数，$t$ 值）}} \\
tone\_P $\times$ 2014后（一日） & $-75.6^{**}$ & $-73.1^{**}$ & $-60.5$ & $-59.0$ \\
        & $(-2.11)$ & $(-2.01)$ & $(-1.25)$ & $(-1.62)$ \\
tone\_P $\times$ 2014后（五日） & $-181.4^{***}$ & $-183.6^{***}$ & $-109.3$ & $-152.5^{**}$ \\
        & $(-2.65)$ & $(-2.71)$ & $(-1.25)$ & $(-2.29)$ \\
\midrule
\multicolumn{5}{l}{\emph{面板 C：置换检验 $p$ 值}} \\
 & 0.020 & — & — & — \\
\bottomrule
\end{tabular}
\begin{tablenotes}[flushleft]
\footnotesize
\item ``替代词典''指以自建货币政策领域词典重新计算情绪指标的结果；``剔除 2015''指去掉 2015 年
四期报告后重估。面板 C 仅报告日内振幅方程的置换 $p$ 值。
\end{tablenotes}
\end{threeparttable}
\end{table}

\begin{table}[htbp]
\centering
\small
\begin{threeparttable}
\caption{补充检验：危机期、安慰剂窗口与政策操作}
\label{tab:robust2}
\begin{tabular}{lcc}
\toprule
设定 & 日内振幅（系数，$t$ 值） & 断点交互项（系数，$t$ 值） \\
\midrule
基准 & $-27.8^{***}$ $(-2.58)$ & $-75.5^{**}$ $(-2.11)$ \\
剔除 2007--2010 & $-15.5^{**}$ $(-2.07)$ & $-36.6$ $(-1.27)$ \\
剔除 2008--2009 & $-18.8^{**}$ $(-2.34)$ & $-45.5$ $(-1.56)$ \\
加入年份趋势 & $-20.7^{**}$ $(-2.18)$ & — \\
事件日前移 20 个交易日 & $+8.5$ $(0.46)$ & $+9.6$ $(0.40)$ \\
事件日后移 20 个交易日 & $-12.5$ $(-1.21)$ & $-35.1$ $(-1.51)$ \\
控制三类政策操作 & $-27.3^{**}$ $(-2.53)$ & $-73.9^{**}$ $(-2.12)$ \\
\bottomrule
\end{tabular}
\begin{tablenotes}[flushleft]
\footnotesize
\item 系数单位为基点，括号内为 Newey--West 稳健 $t$ 值。安慰剂检验按平移后的日期
重新计算振幅与收益，情绪变量不变。政策操作指全面降准、1 年期 LPR 与 7 天逆回购利率，
按报告期汇总为季度强度。``—''表示该项未单独估计。
\end{tablenotes}
\end{threeparttable}
\end{table}

综合来看，本章的证据支持一个有限但明确的判断：政策指引段的文本情绪确实影响市场，
但其作用首先体现在波动而非收益方向上；方向性效应在 2014 年之后才出现，且稳健性
弱于波动结果。这一区分对理解央行沟通的作用机制是重要的——若只报告全样本的平均效应，
两个方向相反的结果会相互抵消，从而得出``沟通无效''的错误结论。
 引入；若主文档是 ctexbook/report，把下面的
%   \section{实证结果} 改成 \chapter{实证结果}，并把 \subsection 降一级。
%   依赖宏包：booktabs、threeparttable、amsmath、graphicx
%   表格编号按"第 5 章"手工写定，若章节顺序调整需同步修改。
% =====================================================================

\section{实证结果}

本章报告文本情绪对市场收益与波动的影响。基准样本为 2001 年第一季度至
2026 年第二季度共 102 期《货币政策执行报告》，其中 99 期可完整切分出政策指引段
与宏观回顾段的情绪指标。市场收益由 16 个连续可得的申万一级行业等权组合构造，
并以 28 个行业口径作为稳健性对照。情绪变量均经标准化处理，回归系数因此可解释为
情绪每变动一个标准差的效应；标准误采用 Newey--West 异方差自相关稳健估计，
滞后阶数取 4。

\subsection{描述性统计}

表~\ref{tab:desc} 列出主要变量的分布。政策指引段情绪（tone\_P）与宏观回顾段情绪
（tone\_R）的相关系数为 0.375，二者相关但并非同一信息的重复度量，这为后续的
分段检验提供了必要的前提。事件日市场收益的均值为 19 个基点，标准差 186 个基点；
日内振幅均值为 2.09\%，意味着报告发布当日最高价与最低价之差平均约为前收盘价的
2\%。事前波动（发布前 60 个交易日市场收益的标准差）均值为 1.46\%。

\begin{table}[htbp]
\centering
\small
\begin{threeparttable}
\caption{描述性统计}
\label{tab:desc}
\begin{tabular}{lcccccc}
\toprule
变量 & 均值 & 标准差 & 最小值 & 中位数 & 最大值 & 观测数 \\
\midrule
政策指引段情绪（tone\_P，标准化） & 0.000 & 1.005 & $-2.773$ & 0.029 & 2.558 & 99 \\
宏观回顾段情绪（tone\_R，标准化） & 0.000 & 1.005 & $-2.928$ & 0.068 & 2.535 & 102 \\
发布后一日收益（\%）              & 0.19  & 1.86  & $-7.60$ & 0.30  & 4.63  & 102 \\
发布后五日收益（\%）              & 0.52  & 3.59  & $-12.59$ & 0.71 & 11.58 & 102 \\
事件日日内振幅（\%）              & 2.09  & 0.92  & 0.80  & 1.88  & 4.76  & 101 \\
事件日绝对收益（\%）              & 0.97  & 0.92  & 0.00  & 0.66  & 3.70  & 102 \\
事前 60 日波动（\%）              & 1.46  & 0.64  & 0.61  & 1.33  & 3.51  & 102 \\
\bottomrule
\end{tabular}
\begin{tablenotes}[flushleft]
\footnotesize
\item 样本为 2001Q1--2026Q2 共 102 期报告；收益与振幅由 16 个申万一级行业的等权组合计算。
\end{tablenotes}
\end{threeparttable}
\end{table}

值得先说明的是情绪与市场波动状态的原始相关性。tone\_P 与事件日日内振幅的相关系数为
$-0.302$，与事件日绝对收益为 $-0.254$，与事前波动的相关系数为 $-0.232$。最后一个
数字提示，情绪较高的报告往往出现在市场相对平静的时期，若不加以控制，很容易把
``情绪降低波动''与``情绪在高波动期偏低''混为一谈。因此在回归中同时报告含事前波动
控制的设定，并以此作为主要判断依据。

\subsection{政策指引情绪与市场波动}

表~\ref{tab:vol} 是本章的主结果。第 (1) 列以事件日绝对收益衡量波动，tone\_P 的系数为
$-23.4$ 个基点（$t=-2.53$）；加入事前波动后（第 (2) 列），系数降至 $-16.6$，显著性
减弱至 10\% 水平，说明有一部分关联确实来自波动状态的持续性。由于绝对收益无法区分
``价格单边移动''与``日内大幅震荡''，第 (3)、(4) 列改用日内振幅，系数分别为 $-27.8$
（$t=-2.58$）与 $-18.5$（$t=-2.43$），控制事前波动后依然在 5\% 水平显著。

第 (5) 列同时纳入回顾段情绪。tone\_R 的系数为 $-2.1$（$t=-0.20$），几乎为零，而
tone\_P 维持在 $-17.8$（$t=-2.27$）。这一对比有明确含义：显著影响波动的并不是
``整篇报告听起来是否乐观''，而是其中面向未来的政策指引部分。第 (6) 列把振幅扩展到
发布后五个交易日，系数为 $-12.1$（$t=-2.03$），效应并未随窗口延长而累积，更接近于
一次性的信息释放而非持续影响。

就经济意义而言，以振幅均值 2.09\% 为参照，$-18.5$ 个基点相当于振幅下降约 9\%；
以标准差衡量，一个标准差的情绪变动对应约 0.2 个标准差的振幅变化。这一量级不算大，
但在 99 个事件、日频噪声显著的条件下已属不易识别。

需要提醒的是，这一幅度对是否包含危机年份较为敏感。剔除 2007--2010 年后，系数降至
$-15.5$（$t=-2.07$），约为基准的一半（见 5.5 节）。方向与显著性不变，但把 $-28$
读作一个稳定不变的效应量并不恰当，更稳妥的表述是：效应量在 15 至 28 个基点之间，
取决于样本是否包含金融危机时期。

分组视角补充了一个细节。把情绪按三分位分组后，低情绪组的平均振幅为 2.53\%，中、高
两组分别为 1.75\% 与 1.99\%（图 5.2）。若把样本切成五组，均值依次为
2.83\%、2.07\%、1.65\%、2.13\%、1.78\%，同样不是单调下降。也就是说，效应主要来自
情绪偏低的一端：措辞异常谨慎的时期波动显著更高，而中高情绪之间并无梯度。这与
``沟通收敛分歧''的解释相容——市场在意的未必是措辞有多乐观，而是它是否透露出不安。

\begin{table}[htbp]
\centering
\small
\begin{threeparttable}
\caption{政策指引段情绪与市场波动}
\label{tab:vol}
\begin{tabular}{lcccccc}
\toprule
 & (1) & (2) & (3) & (4) & (5) & (6) \\
 & $|r|$ & $|r|$ & 振幅 & 振幅 & 振幅 & 五日振幅 \\
\midrule
tone\_P        & $-23.4^{**}$ & $-16.6^{*}$ & $-27.8^{***}$ & $-18.5^{**}$ & $-17.8^{**}$ & $-12.1^{**}$ \\
                & $(-2.53)$ & $(-1.96)$ & $(-2.58)$ & $(-2.43)$ & $(-2.27)$ & $(-2.03)$ \\
tone\_R        & & & & & $-2.1$ & \\
                & & & & & $(-0.20)$ & \\
事前波动        & & 46.0$^{***}$ & & 61.9$^{***}$ & 61.1$^{***}$ & 80.0$^{***}$ \\
                & & $(4.10)$ & & $(4.24)$ & $(4.14)$ & $(4.16)$ \\
\midrule
观测数          & 99 & 99 & 98 & 98 & 98 & 99 \\
$R^2$           & 0.065 & 0.162 & 0.091 & 0.267 & 0.268 & 0.355 \\
\bottomrule
\end{tabular}
\begin{tablenotes}[flushleft]
\footnotesize
\item 被解释变量为事件日市场波动，单位为百分比；情绪与事前波动的系数均已换算为基点
（1 个基点 $=0.01$ 个百分点），其中事前波动的系数读作``事前波动每上升 1 个百分点''。
括号内为 Newey--West 稳健 $t$ 值。$^{*}$、$^{**}$、$^{***}$ 分别表示 10\%、5\%、1\% 显著性水平。
\end{tablenotes}
\end{threeparttable}
\end{table}

\subsection{制度断点：2014 年前后}

样本期内货币政策框架发生了明显变化：2013 年流动性紧张、2014 年利率市场化加速、
2015 年汇率形成机制调整，报告自身的章节结构也在 2015 年前后改写。这些变化提示
情绪的作用方式未必稳定。表~\ref{tab:break} 在方程中引入 2014 年虚拟变量及其与
情绪的交互项，直接检验这一可能。

收益方程中，交互项系数为 $-75.6$ 个基点（$t=-2.11$），五日窗口扩大到 $-181.4$
（$t=-2.65$），即 2014 年以前情绪对收益方向无显著影响（点估计 $+31.6$，$t=0.97$），
2014 年以后转为负向。波动方程则呈现相反的变化：2014 年以前的系数为 $-35.0$
（$t=-3.10$），交互项为 $+26.6$（$t=1.59$），说明波动渠道在后期减弱，但未完全消失。

需要说明的是断点的选择并非由数据挑选而来，而是对应上述制度节点。为检验这一处理是否
稳健，表~\ref{tab:breaksens} 将断点依次替换为 2008、2010、2012、2014、2016、2018 年。
2008 年作断点时交互项接近于零，其余各年份的系数方向一致，其中 2010、2012、2014、2018
年显著，2016 年偏弱。这说明变化是渐进的，把它归因于某一个特定年份并不合适，但
``2010 年代之后效应改变''这一判断是站得住的。

\begin{table}[htbp]
\centering
\small
\begin{threeparttable}
\caption{制度断点检验}
\label{tab:break}
\begin{tabular}{lcccc}
\toprule
 & (1) & (2) & (3) & (4) \\
 & 一日收益 & 五日收益 & 日内振幅 & $|r|$ \\
\midrule
tone\_P              & $+31.6$   & $+73.7$   & $-35.0^{***}$ & $-32.7^{***}$ \\
                      & $(0.97)$  & $(1.22)$  & $(-3.10)$     & $(-2.97)$ \\
tone\_P $\times$ 2014后 & $-75.6^{**}$ & $-181.4^{***}$ & $+26.6$ & $+25.1$ \\
                      & $(-2.11)$ & $(-2.65)$ & $(1.59)$      & $(1.53)$ \\
\midrule
观测数                & 99 & 99 & 98 & 99 \\
$R^2$                 & 0.076 & 0.077 & 0.190 & 0.084 \\
\bottomrule
\end{tabular}
\begin{tablenotes}[flushleft]
\footnotesize
\item 各列均控制 2014 年虚拟变量本身；系数单位为基点，括号内为 Newey--West 稳健 $t$ 值。
\end{tablenotes}
\end{threeparttable}
\end{table}

\begin{table}[htbp]
\centering
\small
\begin{threeparttable}
\caption{断点年份敏感性（交互项系数与 $t$ 值）}
\label{tab:breaksens}
\begin{tabular}{lccc}
\toprule
断点年份 & 一日收益 & 五日收益 & 日内振幅 \\
\midrule
2008 & $-5.9\ (-0.14)$   & $-1.6\ (-0.02)$    & $+1.8\ (0.10)$ \\
2010 & $-86.6^{**}\ (-2.41)$ & $-193.1^{***}\ (-2.69)$ & $+34.5^{**}\ (2.04)$ \\
2012 & $-76.6^{**}\ (-2.20)$ & $-183.5^{***}\ (-2.74)$ & $+27.4^{*}\ (1.71)$ \\
2014 & $-75.6^{**}\ (-2.11)$ & $-181.4^{***}\ (-2.65)$ & $+26.6\ (1.59)$ \\
2016 & $-57.7\ (-1.47)$  & $-134.9^{*}\ (-1.88)$   & $+34.4^{**}\ (2.07)$ \\
2018 & $-64.4^{*}\ (-1.82)$ & $-164.9^{**}\ (-2.29)$ & $+16.2\ (0.79)$ \\
\bottomrule
\end{tabular}
\begin{tablenotes}[flushleft]
\footnotesize
\item 每行对应一次单独的回归，仅改变断点年份，其余设定与表~\ref{tab:break} 相同。
\end{tablenotes}
\end{threeparttable}
\end{table}

\subsection{2014 年之后的方向性效应}

既然断点检验显示后期情绪与收益转为负相关，这里单独报告 2014 年以后的子样本。
表~\ref{tab:post14} 显示，tone\_P 对发布后一日收益的系数为 $-44.0$ 个基点
（$t=-2.89$），五日窗口为 $-107.8$（$t=-3.30$）；同时，波动类方程在子样本中
全部不显著。两组结果互为印证：后期的定价体现在方向上，而不在波动上。

对这一结果需要保持克制。首先，负号的含义并非``情绪越宽松市场越差''这么简单，
更可能反映宽松措辞集中出现在基本面走弱的时点，市场将政策宽松解读为经济前景的
负面信号。其次，该结果对危机年份高度敏感：剔除 2015 年后，一日窗口的交互项由
$-75.6$（$t=-2.11$）降至 $-59.0$（$t=-1.62$）；剔除 2008--2009 年后降至 $-45.5$
（$t=-1.56$）。任一段被剔除，5\% 水平的显著性即不复存在。因此本文不把它作为独立
结论，仅记录符号的方向性变化。

\begin{table}[htbp]
\centering
\small
\begin{threeparttable}
\caption{2014 年以后的子样本}
\label{tab:post14}
\begin{tabular}{lcccc}
\toprule
 & (1) & (2) & (3) & (4) \\
 & 一日收益 & 五日收益 & 一日·双段 & 日内振幅 \\
\midrule
tone\_P & $-44.0^{***}$ & $-107.8^{***}$ & $-47.2^{**}$ & $-8.4$ \\
        & $(-2.89)$     & $(-3.30)$      & $(-2.36)$    & $(-0.68)$ \\
tone\_R &               &                & $+5.9$       & \\
        &               &                & $(0.28)$     & \\
\midrule
观测数  & 50 & 50 & 50 & 49 \\
\bottomrule
\end{tabular}
\begin{tablenotes}[flushleft]
\footnotesize
\item 子样本为 2014Q1--2026Q2。系数单位为基点，括号内为 Newey--West 稳健 $t$ 值。
\end{tablenotes}
\end{threeparttable}
\end{table}

\subsection{稳健性检验}

表~\ref{tab:robust} 汇总四个方面的检验，表~\ref{tab:robust2} 补充三项针对性检验。

\emph{市场口径}。把等权组合由 16 个行业扩展为 28 个行业，主结果基本不变：
日内振幅系数 $-26.7$（$t=-2.58$），断点交互项 $-73.1$（$t=-2.01$），五日交互项
$-183.6$（$t=-2.71$）。结论不依赖于行业集合的选择。

\emph{词典}。以本文自建的货币政策领域词典（积极 159 词、消极 139 词）重新度量情绪，
各系数方向与主结果一致，但量级缩小、显著性下降：绝对收益方程仍显著（$-19.1$，
$t=-2.43$），日内振幅方程降至 $-7.8$（$t=-0.66$），断点交互项降至 $-60.5$
（$t=-1.25$）。该词典每千字仅命中约 68 个情感词，主词典为 104 个，覆盖率差异恰好
对应系数的等比例收缩，符合测量误差导致衰减的预期。因此这一结果是``方向稳健、精度不足''，
而非``结论被推翻''。

\emph{影响点}。逐一剔除 2015 年（股市异常波动）与 2020 年（疫情冲击），并对收益做
1\%/99\% 缩尾。波动类结果对上述处理均不敏感，日内振幅系数稳定在 $-25.8$ 至 $-27.9$
之间。断点交互项则较为敏感：剔除 2015 年后，一日窗口的交互项由 $-75.6$（$t=-2.11$）
降至 $-59.0$（$t=-1.62$），显著性消失；五日窗口降至 $-152.5$（$t=-2.29$），仍在 5\%
水平显著。剔除 2020 年与收益缩尾后，一日交互项分别为 $-73.0$ 与 $-74.4$，保持在 5\%
水平。可见后期方向效应确有相当一部分由 2015 年的数期报告贡献，这解释了上文为何
采取保留态度，也说明该效应在更长窗口上比在单日窗口上更稳定。

\emph{置换检验}。将情绪序列与事件随机配对 500 次，观察真实 $t$ 值在随机分布中的位置。
日内振幅方程的置换 $p$ 值为 0.020，加入回顾段与事前波动后为 0.068，断点交互项为 0.058。
置换检验不依赖分布假设，其结果与稳健 $t$ 检验基本一致：波动结果在 5\% 水平成立，
断点结果处于边际。

\emph{危机期敏感性}。2007--2009 年的措辞与波动同时处于极端，是本文最担心的混淆来源。
剔除 2007--2010 年后，振幅方程系数降至 $-15.5$（$t=-2.07$），剔除 2008--2009 年后为
$-18.8$（$t=-2.34$），加入年份趋势后为 $-20.7$（$t=-2.18$）。方向与显著性保持不变，
幅度约为基准的一半。断点交互项则不同：剔除 2008--2009 年后降至 $-45.5$（$t=-1.56$），
不再显著。可见波动结果稳健但幅度受危机年份影响，方向结果则更依赖少数极端时期。

\emph{安慰剂窗口}。把事件日整体平移 $\pm$5、$\pm$10、$\pm$20 个交易日重估。平移 $\pm$10
天以上时，振幅方程的系数不再显著（$t$ 值均小于 1.1），断点交互项也不再成立，说明效应
与发布时点相关。但平移 $-5$ 天时系数仍有 $-26.6$（$t=-1.79$），保留了约九成的幅度。
这表明效应并非严格的单日冲击，而是集中在发布前后约一周的窗口内，与报告部分内容提前
释放的解释相容。

\emph{政策操作控制}。一个自然的担心是，文本情绪只是货币政策操作的代理。本文整理了三类
总量操作——全面降准 24 次、1 年期 LPR 下调 12 次、7 天逆回购利率下调 12 次——按报告期
汇总为季度操作强度，加入振幅方程。情绪系数由 $-27.8$ 变为 $-27.3$（$t=-2.53$），加入
事前波动后为 $-16.7$（$t=-2.17$）；三类操作本身均不显著，与情绪的相关性也只在 $-0.01$
到 $-0.21$ 之间。断点交互项由 $-75.5$ 变为 $-73.9$（$t=-2.12$）。这说明波动效应并非
由降准降息驱动。

\begin{table}[htbp]
\centering
\small
\begin{threeparttable}
\caption{稳健性检验汇总}
\label{tab:robust}
\begin{tabular}{lcccc}
\toprule
 & 主结果 & 28 行业 & 替代词典 & 剔除 2015 \\
\midrule
\multicolumn{5}{l}{\emph{面板 A：日内振幅（系数，$t$ 值）}} \\
tone\_P & $-27.8^{***}$ & $-26.7^{***}$ & $-7.8$ & $-25.8^{**}$ \\
        & $(-2.58)$ & $(-2.58)$ & $(-0.66)$ & $(-2.19)$ \\
\midrule
\multicolumn{5}{l}{\emph{面板 B：断点交互项（系数，$t$ 值）}} \\
tone\_P $\times$ 2014后（一日） & $-75.6^{**}$ & $-73.1^{**}$ & $-60.5$ & $-59.0$ \\
        & $(-2.11)$ & $(-2.01)$ & $(-1.25)$ & $(-1.62)$ \\
tone\_P $\times$ 2014后（五日） & $-181.4^{***}$ & $-183.6^{***}$ & $-109.3$ & $-152.5^{**}$ \\
        & $(-2.65)$ & $(-2.71)$ & $(-1.25)$ & $(-2.29)$ \\
\midrule
\multicolumn{5}{l}{\emph{面板 C：置换检验 $p$ 值}} \\
 & 0.020 & — & — & — \\
\bottomrule
\end{tabular}
\begin{tablenotes}[flushleft]
\footnotesize
\item ``替代词典''指以自建货币政策领域词典重新计算情绪指标的结果；``剔除 2015''指去掉 2015 年
四期报告后重估。面板 C 仅报告日内振幅方程的置换 $p$ 值。
\end{tablenotes}
\end{threeparttable}
\end{table}

\begin{table}[htbp]
\centering
\small
\begin{threeparttable}
\caption{补充检验：危机期、安慰剂窗口与政策操作}
\label{tab:robust2}
\begin{tabular}{lcc}
\toprule
设定 & 日内振幅（系数，$t$ 值） & 断点交互项（系数，$t$ 值） \\
\midrule
基准 & $-27.8^{***}$ $(-2.58)$ & $-75.5^{**}$ $(-2.11)$ \\
剔除 2007--2010 & $-15.5^{**}$ $(-2.07)$ & $-36.6$ $(-1.27)$ \\
剔除 2008--2009 & $-18.8^{**}$ $(-2.34)$ & $-45.5$ $(-1.56)$ \\
加入年份趋势 & $-20.7^{**}$ $(-2.18)$ & — \\
事件日前移 20 个交易日 & $+8.5$ $(0.46)$ & $+9.6$ $(0.40)$ \\
事件日后移 20 个交易日 & $-12.5$ $(-1.21)$ & $-35.1$ $(-1.51)$ \\
控制三类政策操作 & $-27.3^{**}$ $(-2.53)$ & $-73.9^{**}$ $(-2.12)$ \\
\bottomrule
\end{tabular}
\begin{tablenotes}[flushleft]
\footnotesize
\item 系数单位为基点，括号内为 Newey--West 稳健 $t$ 值。安慰剂检验按平移后的日期
重新计算振幅与收益，情绪变量不变。政策操作指全面降准、1 年期 LPR 与 7 天逆回购利率，
按报告期汇总为季度强度。``—''表示该项未单独估计。
\end{tablenotes}
\end{threeparttable}
\end{table}

综合来看，本章的证据支持一个有限但明确的判断：政策指引段的文本情绪确实影响市场，
但其作用首先体现在波动而非收益方向上；方向性效应在 2014 年之后才出现，且稳健性
弱于波动结果。这一区分对理解央行沟通的作用机制是重要的——若只报告全样本的平均效应，
两个方向相反的结果会相互抵消，从而得出``沟通无效''的错误结论。
 引入；若主文档是 ctexbook/report，把下面的
%   \section{实证结果} 改成 \chapter{实证结果}，并把 \subsection 降一级。
%   依赖宏包：booktabs、threeparttable、amsmath、graphicx
%   表格编号按"第 5 章"手工写定，若章节顺序调整需同步修改。
% =====================================================================

\section{实证结果}

本章报告文本情绪对市场收益与波动的影响。基准样本为 2001 年第一季度至
2026 年第二季度共 102 期《货币政策执行报告》，其中 99 期可完整切分出政策指引段
与宏观回顾段的情绪指标。市场收益由 16 个连续可得的申万一级行业等权组合构造，
并以 28 个行业口径作为稳健性对照。情绪变量均经标准化处理，回归系数因此可解释为
情绪每变动一个标准差的效应；标准误采用 Newey--West 异方差自相关稳健估计，
滞后阶数取 4。

\subsection{描述性统计}

表~\ref{tab:desc} 列出主要变量的分布。政策指引段情绪（tone\_P）与宏观回顾段情绪
（tone\_R）的相关系数为 0.375，二者相关但并非同一信息的重复度量，这为后续的
分段检验提供了必要的前提。事件日市场收益的均值为 19 个基点，标准差 186 个基点；
日内振幅均值为 2.09\%，意味着报告发布当日最高价与最低价之差平均约为前收盘价的
2\%。事前波动（发布前 60 个交易日市场收益的标准差）均值为 1.46\%。

\begin{table}[htbp]
\centering
\small
\begin{threeparttable}
\caption{描述性统计}
\label{tab:desc}
\begin{tabular}{lcccccc}
\toprule
变量 & 均值 & 标准差 & 最小值 & 中位数 & 最大值 & 观测数 \\
\midrule
政策指引段情绪（tone\_P，标准化） & 0.000 & 1.005 & $-2.773$ & 0.029 & 2.558 & 99 \\
宏观回顾段情绪（tone\_R，标准化） & 0.000 & 1.005 & $-2.928$ & 0.068 & 2.535 & 102 \\
发布后一日收益（\%）              & 0.19  & 1.86  & $-7.60$ & 0.30  & 4.63  & 102 \\
发布后五日收益（\%）              & 0.52  & 3.59  & $-12.59$ & 0.71 & 11.58 & 102 \\
事件日日内振幅（\%）              & 2.09  & 0.92  & 0.80  & 1.88  & 4.76  & 101 \\
事件日绝对收益（\%）              & 0.97  & 0.92  & 0.00  & 0.66  & 3.70  & 102 \\
事前 60 日波动（\%）              & 1.46  & 0.64  & 0.61  & 1.33  & 3.51  & 102 \\
\bottomrule
\end{tabular}
\begin{tablenotes}[flushleft]
\footnotesize
\item 样本为 2001Q1--2026Q2 共 102 期报告；收益与振幅由 16 个申万一级行业的等权组合计算。
\end{tablenotes}
\end{threeparttable}
\end{table}

值得先说明的是情绪与市场波动状态的原始相关性。tone\_P 与事件日日内振幅的相关系数为
$-0.302$，与事件日绝对收益为 $-0.254$，与事前波动的相关系数为 $-0.232$。最后一个
数字提示，情绪较高的报告往往出现在市场相对平静的时期，若不加以控制，很容易把
``情绪降低波动''与``情绪在高波动期偏低''混为一谈。因此在回归中同时报告含事前波动
控制的设定，并以此作为主要判断依据。

\subsection{政策指引情绪与市场波动}

表~\ref{tab:vol} 是本章的主结果。第 (1) 列以事件日绝对收益衡量波动，tone\_P 的系数为
$-23.4$ 个基点（$t=-2.53$）；加入事前波动后（第 (2) 列），系数降至 $-16.6$，显著性
减弱至 10\% 水平，说明有一部分关联确实来自波动状态的持续性。由于绝对收益无法区分
``价格单边移动''与``日内大幅震荡''，第 (3)、(4) 列改用日内振幅，系数分别为 $-27.8$
（$t=-2.58$）与 $-18.5$（$t=-2.43$），控制事前波动后依然在 5\% 水平显著。

第 (5) 列同时纳入回顾段情绪。tone\_R 的系数为 $-2.1$（$t=-0.20$），几乎为零，而
tone\_P 维持在 $-17.8$（$t=-2.27$）。这一对比有明确含义：显著影响波动的并不是
``整篇报告听起来是否乐观''，而是其中面向未来的政策指引部分。第 (6) 列把振幅扩展到
发布后五个交易日，系数为 $-12.1$（$t=-2.03$），效应并未随窗口延长而累积，更接近于
一次性的信息释放而非持续影响。

就经济意义而言，以振幅均值 2.09\% 为参照，$-18.5$ 个基点相当于振幅下降约 9\%；
以标准差衡量，一个标准差的情绪变动对应约 0.2 个标准差的振幅变化。这一量级不算大，
但在 99 个事件、日频噪声显著的条件下已属不易识别。

需要提醒的是，这一幅度对是否包含危机年份较为敏感。剔除 2007--2010 年后，系数降至
$-15.5$（$t=-2.07$），约为基准的一半（见 5.5 节）。方向与显著性不变，但把 $-28$
读作一个稳定不变的效应量并不恰当，更稳妥的表述是：效应量在 15 至 28 个基点之间，
取决于样本是否包含金融危机时期。

分组视角补充了一个细节。把情绪按三分位分组后，低情绪组的平均振幅为 2.53\%，中、高
两组分别为 1.75\% 与 1.99\%（图 5.2）。若把样本切成五组，均值依次为
2.83\%、2.07\%、1.65\%、2.13\%、1.78\%，同样不是单调下降。也就是说，效应主要来自
情绪偏低的一端：措辞异常谨慎的时期波动显著更高，而中高情绪之间并无梯度。这与
``沟通收敛分歧''的解释相容——市场在意的未必是措辞有多乐观，而是它是否透露出不安。

\begin{table}[htbp]
\centering
\small
\begin{threeparttable}
\caption{政策指引段情绪与市场波动}
\label{tab:vol}
\begin{tabular}{lcccccc}
\toprule
 & (1) & (2) & (3) & (4) & (5) & (6) \\
 & $|r|$ & $|r|$ & 振幅 & 振幅 & 振幅 & 五日振幅 \\
\midrule
tone\_P        & $-23.4^{**}$ & $-16.6^{*}$ & $-27.8^{***}$ & $-18.5^{**}$ & $-17.8^{**}$ & $-12.1^{**}$ \\
                & $(-2.53)$ & $(-1.96)$ & $(-2.58)$ & $(-2.43)$ & $(-2.27)$ & $(-2.03)$ \\
tone\_R        & & & & & $-2.1$ & \\
                & & & & & $(-0.20)$ & \\
事前波动        & & 46.0$^{***}$ & & 61.9$^{***}$ & 61.1$^{***}$ & 80.0$^{***}$ \\
                & & $(4.10)$ & & $(4.24)$ & $(4.14)$ & $(4.16)$ \\
\midrule
观测数          & 99 & 99 & 98 & 98 & 98 & 99 \\
$R^2$           & 0.065 & 0.162 & 0.091 & 0.267 & 0.268 & 0.355 \\
\bottomrule
\end{tabular}
\begin{tablenotes}[flushleft]
\footnotesize
\item 被解释变量为事件日市场波动，单位为百分比；情绪与事前波动的系数均已换算为基点
（1 个基点 $=0.01$ 个百分点），其中事前波动的系数读作``事前波动每上升 1 个百分点''。
括号内为 Newey--West 稳健 $t$ 值。$^{*}$、$^{**}$、$^{***}$ 分别表示 10\%、5\%、1\% 显著性水平。
\end{tablenotes}
\end{threeparttable}
\end{table}

\subsection{制度断点：2014 年前后}

样本期内货币政策框架发生了明显变化：2013 年流动性紧张、2014 年利率市场化加速、
2015 年汇率形成机制调整，报告自身的章节结构也在 2015 年前后改写。这些变化提示
情绪的作用方式未必稳定。表~\ref{tab:break} 在方程中引入 2014 年虚拟变量及其与
情绪的交互项，直接检验这一可能。

收益方程中，交互项系数为 $-75.6$ 个基点（$t=-2.11$），五日窗口扩大到 $-181.4$
（$t=-2.65$），即 2014 年以前情绪对收益方向无显著影响（点估计 $+31.6$，$t=0.97$），
2014 年以后转为负向。波动方程则呈现相反的变化：2014 年以前的系数为 $-35.0$
（$t=-3.10$），交互项为 $+26.6$（$t=1.59$），说明波动渠道在后期减弱，但未完全消失。

需要说明的是断点的选择并非由数据挑选而来，而是对应上述制度节点。为检验这一处理是否
稳健，表~\ref{tab:breaksens} 将断点依次替换为 2008、2010、2012、2014、2016、2018 年。
2008 年作断点时交互项接近于零，其余各年份的系数方向一致，其中 2010、2012、2014、2018
年显著，2016 年偏弱。这说明变化是渐进的，把它归因于某一个特定年份并不合适，但
``2010 年代之后效应改变''这一判断是站得住的。

\begin{table}[htbp]
\centering
\small
\begin{threeparttable}
\caption{制度断点检验}
\label{tab:break}
\begin{tabular}{lcccc}
\toprule
 & (1) & (2) & (3) & (4) \\
 & 一日收益 & 五日收益 & 日内振幅 & $|r|$ \\
\midrule
tone\_P              & $+31.6$   & $+73.7$   & $-35.0^{***}$ & $-32.7^{***}$ \\
                      & $(0.97)$  & $(1.22)$  & $(-3.10)$     & $(-2.97)$ \\
tone\_P $\times$ 2014后 & $-75.6^{**}$ & $-181.4^{***}$ & $+26.6$ & $+25.1$ \\
                      & $(-2.11)$ & $(-2.65)$ & $(1.59)$      & $(1.53)$ \\
\midrule
观测数                & 99 & 99 & 98 & 99 \\
$R^2$                 & 0.076 & 0.077 & 0.190 & 0.084 \\
\bottomrule
\end{tabular}
\begin{tablenotes}[flushleft]
\footnotesize
\item 各列均控制 2014 年虚拟变量本身；系数单位为基点，括号内为 Newey--West 稳健 $t$ 值。
\end{tablenotes}
\end{threeparttable}
\end{table}

\begin{table}[htbp]
\centering
\small
\begin{threeparttable}
\caption{断点年份敏感性（交互项系数与 $t$ 值）}
\label{tab:breaksens}
\begin{tabular}{lccc}
\toprule
断点年份 & 一日收益 & 五日收益 & 日内振幅 \\
\midrule
2008 & $-5.9\ (-0.14)$   & $-1.6\ (-0.02)$    & $+1.8\ (0.10)$ \\
2010 & $-86.6^{**}\ (-2.41)$ & $-193.1^{***}\ (-2.69)$ & $+34.5^{**}\ (2.04)$ \\
2012 & $-76.6^{**}\ (-2.20)$ & $-183.5^{***}\ (-2.74)$ & $+27.4^{*}\ (1.71)$ \\
2014 & $-75.6^{**}\ (-2.11)$ & $-181.4^{***}\ (-2.65)$ & $+26.6\ (1.59)$ \\
2016 & $-57.7\ (-1.47)$  & $-134.9^{*}\ (-1.88)$   & $+34.4^{**}\ (2.07)$ \\
2018 & $-64.4^{*}\ (-1.82)$ & $-164.9^{**}\ (-2.29)$ & $+16.2\ (0.79)$ \\
\bottomrule
\end{tabular}
\begin{tablenotes}[flushleft]
\footnotesize
\item 每行对应一次单独的回归，仅改变断点年份，其余设定与表~\ref{tab:break} 相同。
\end{tablenotes}
\end{threeparttable}
\end{table}

\subsection{2014 年之后的方向性效应}

既然断点检验显示后期情绪与收益转为负相关，这里单独报告 2014 年以后的子样本。
表~\ref{tab:post14} 显示，tone\_P 对发布后一日收益的系数为 $-44.0$ 个基点
（$t=-2.89$），五日窗口为 $-107.8$（$t=-3.30$）；同时，波动类方程在子样本中
全部不显著。两组结果互为印证：后期的定价体现在方向上，而不在波动上。

对这一结果需要保持克制。首先，负号的含义并非``情绪越宽松市场越差''这么简单，
更可能反映宽松措辞集中出现在基本面走弱的时点，市场将政策宽松解读为经济前景的
负面信号。其次，该结果对危机年份高度敏感：剔除 2015 年后，一日窗口的交互项由
$-75.6$（$t=-2.11$）降至 $-59.0$（$t=-1.62$）；剔除 2008--2009 年后降至 $-45.5$
（$t=-1.56$）。任一段被剔除，5\% 水平的显著性即不复存在。因此本文不把它作为独立
结论，仅记录符号的方向性变化。

\begin{table}[htbp]
\centering
\small
\begin{threeparttable}
\caption{2014 年以后的子样本}
\label{tab:post14}
\begin{tabular}{lcccc}
\toprule
 & (1) & (2) & (3) & (4) \\
 & 一日收益 & 五日收益 & 一日·双段 & 日内振幅 \\
\midrule
tone\_P & $-44.0^{***}$ & $-107.8^{***}$ & $-47.2^{**}$ & $-8.4$ \\
        & $(-2.89)$     & $(-3.30)$      & $(-2.36)$    & $(-0.68)$ \\
tone\_R &               &                & $+5.9$       & \\
        &               &                & $(0.28)$     & \\
\midrule
观测数  & 50 & 50 & 50 & 49 \\
\bottomrule
\end{tabular}
\begin{tablenotes}[flushleft]
\footnotesize
\item 子样本为 2014Q1--2026Q2。系数单位为基点，括号内为 Newey--West 稳健 $t$ 值。
\end{tablenotes}
\end{threeparttable}
\end{table}

\subsection{稳健性检验}

表~\ref{tab:robust} 汇总四个方面的检验，表~\ref{tab:robust2} 补充三项针对性检验。

\emph{市场口径}。把等权组合由 16 个行业扩展为 28 个行业，主结果基本不变：
日内振幅系数 $-26.7$（$t=-2.58$），断点交互项 $-73.1$（$t=-2.01$），五日交互项
$-183.6$（$t=-2.71$）。结论不依赖于行业集合的选择。

\emph{词典}。以本文自建的货币政策领域词典（积极 159 词、消极 139 词）重新度量情绪，
各系数方向与主结果一致，但量级缩小、显著性下降：绝对收益方程仍显著（$-19.1$，
$t=-2.43$），日内振幅方程降至 $-7.8$（$t=-0.66$），断点交互项降至 $-60.5$
（$t=-1.25$）。该词典每千字仅命中约 68 个情感词，主词典为 104 个，覆盖率差异恰好
对应系数的等比例收缩，符合测量误差导致衰减的预期。因此这一结果是``方向稳健、精度不足''，
而非``结论被推翻''。

\emph{影响点}。逐一剔除 2015 年（股市异常波动）与 2020 年（疫情冲击），并对收益做
1\%/99\% 缩尾。波动类结果对上述处理均不敏感，日内振幅系数稳定在 $-25.8$ 至 $-27.9$
之间。断点交互项则较为敏感：剔除 2015 年后，一日窗口的交互项由 $-75.6$（$t=-2.11$）
降至 $-59.0$（$t=-1.62$），显著性消失；五日窗口降至 $-152.5$（$t=-2.29$），仍在 5\%
水平显著。剔除 2020 年与收益缩尾后，一日交互项分别为 $-73.0$ 与 $-74.4$，保持在 5\%
水平。可见后期方向效应确有相当一部分由 2015 年的数期报告贡献，这解释了上文为何
采取保留态度，也说明该效应在更长窗口上比在单日窗口上更稳定。

\emph{置换检验}。将情绪序列与事件随机配对 500 次，观察真实 $t$ 值在随机分布中的位置。
日内振幅方程的置换 $p$ 值为 0.020，加入回顾段与事前波动后为 0.068，断点交互项为 0.058。
置换检验不依赖分布假设，其结果与稳健 $t$ 检验基本一致：波动结果在 5\% 水平成立，
断点结果处于边际。

\emph{危机期敏感性}。2007--2009 年的措辞与波动同时处于极端，是本文最担心的混淆来源。
剔除 2007--2010 年后，振幅方程系数降至 $-15.5$（$t=-2.07$），剔除 2008--2009 年后为
$-18.8$（$t=-2.34$），加入年份趋势后为 $-20.7$（$t=-2.18$）。方向与显著性保持不变，
幅度约为基准的一半。断点交互项则不同：剔除 2008--2009 年后降至 $-45.5$（$t=-1.56$），
不再显著。可见波动结果稳健但幅度受危机年份影响，方向结果则更依赖少数极端时期。

\emph{安慰剂窗口}。把事件日整体平移 $\pm$5、$\pm$10、$\pm$20 个交易日重估。平移 $\pm$10
天以上时，振幅方程的系数不再显著（$t$ 值均小于 1.1），断点交互项也不再成立，说明效应
与发布时点相关。但平移 $-5$ 天时系数仍有 $-26.6$（$t=-1.79$），保留了约九成的幅度。
这表明效应并非严格的单日冲击，而是集中在发布前后约一周的窗口内，与报告部分内容提前
释放的解释相容。

\emph{政策操作控制}。一个自然的担心是，文本情绪只是货币政策操作的代理。本文整理了三类
总量操作——全面降准 24 次、1 年期 LPR 下调 12 次、7 天逆回购利率下调 12 次——按报告期
汇总为季度操作强度，加入振幅方程。情绪系数由 $-27.8$ 变为 $-27.3$（$t=-2.53$），加入
事前波动后为 $-16.7$（$t=-2.17$）；三类操作本身均不显著，与情绪的相关性也只在 $-0.01$
到 $-0.21$ 之间。断点交互项由 $-75.5$ 变为 $-73.9$（$t=-2.12$）。这说明波动效应并非
由降准降息驱动。

\begin{table}[htbp]
\centering
\small
\begin{threeparttable}
\caption{稳健性检验汇总}
\label{tab:robust}
\begin{tabular}{lcccc}
\toprule
 & 主结果 & 28 行业 & 替代词典 & 剔除 2015 \\
\midrule
\multicolumn{5}{l}{\emph{面板 A：日内振幅（系数，$t$ 值）}} \\
tone\_P & $-27.8^{***}$ & $-26.7^{***}$ & $-7.8$ & $-25.8^{**}$ \\
        & $(-2.58)$ & $(-2.58)$ & $(-0.66)$ & $(-2.19)$ \\
\midrule
\multicolumn{5}{l}{\emph{面板 B：断点交互项（系数，$t$ 值）}} \\
tone\_P $\times$ 2014后（一日） & $-75.6^{**}$ & $-73.1^{**}$ & $-60.5$ & $-59.0$ \\
        & $(-2.11)$ & $(-2.01)$ & $(-1.25)$ & $(-1.62)$ \\
tone\_P $\times$ 2014后（五日） & $-181.4^{***}$ & $-183.6^{***}$ & $-109.3$ & $-152.5^{**}$ \\
        & $(-2.65)$ & $(-2.71)$ & $(-1.25)$ & $(-2.29)$ \\
\midrule
\multicolumn{5}{l}{\emph{面板 C：置换检验 $p$ 值}} \\
 & 0.020 & — & — & — \\
\bottomrule
\end{tabular}
\begin{tablenotes}[flushleft]
\footnotesize
\item ``替代词典''指以自建货币政策领域词典重新计算情绪指标的结果；``剔除 2015''指去掉 2015 年
四期报告后重估。面板 C 仅报告日内振幅方程的置换 $p$ 值。
\end{tablenotes}
\end{threeparttable}
\end{table}

\begin{table}[htbp]
\centering
\small
\begin{threeparttable}
\caption{补充检验：危机期、安慰剂窗口与政策操作}
\label{tab:robust2}
\begin{tabular}{lcc}
\toprule
设定 & 日内振幅（系数，$t$ 值） & 断点交互项（系数，$t$ 值） \\
\midrule
基准 & $-27.8^{***}$ $(-2.58)$ & $-75.5^{**}$ $(-2.11)$ \\
剔除 2007--2010 & $-15.5^{**}$ $(-2.07)$ & $-36.6$ $(-1.27)$ \\
剔除 2008--2009 & $-18.8^{**}$ $(-2.34)$ & $-45.5$ $(-1.56)$ \\
加入年份趋势 & $-20.7^{**}$ $(-2.18)$ & — \\
事件日前移 20 个交易日 & $+8.5$ $(0.46)$ & $+9.6$ $(0.40)$ \\
事件日后移 20 个交易日 & $-12.5$ $(-1.21)$ & $-35.1$ $(-1.51)$ \\
控制三类政策操作 & $-27.3^{**}$ $(-2.53)$ & $-73.9^{**}$ $(-2.12)$ \\
\bottomrule
\end{tabular}
\begin{tablenotes}[flushleft]
\footnotesize
\item 系数单位为基点，括号内为 Newey--West 稳健 $t$ 值。安慰剂检验按平移后的日期
重新计算振幅与收益，情绪变量不变。政策操作指全面降准、1 年期 LPR 与 7 天逆回购利率，
按报告期汇总为季度强度。``—''表示该项未单独估计。
\end{tablenotes}
\end{threeparttable}
\end{table}

综合来看，本章的证据支持一个有限但明确的判断：政策指引段的文本情绪确实影响市场，
但其作用首先体现在波动而非收益方向上；方向性效应在 2014 年之后才出现，且稳健性
弱于波动结果。这一区分对理解央行沟通的作用机制是重要的——若只报告全样本的平均效应，
两个方向相反的结果会相互抵消，从而得出``沟通无效''的错误结论。
 引入；若主文档是 ctexbook/report，把下面的
%   \section{实证结果} 改成 \chapter{实证结果}，并把 \subsection 降一级。
%   依赖宏包：booktabs、threeparttable、amsmath、graphicx
%   表格编号按"第 5 章"手工写定，若章节顺序调整需同步修改。
% =====================================================================

\section{实证结果}

本章报告文本情绪对市场收益与波动的影响。基准样本为 2001 年第一季度至
2026 年第二季度共 102 期《货币政策执行报告》，其中 99 期可完整切分出政策指引段
与宏观回顾段的情绪指标。市场收益由 16 个连续可得的申万一级行业等权组合构造，
并以 28 个行业口径作为稳健性对照。情绪变量均经标准化处理，回归系数因此可解释为
情绪每变动一个标准差的效应；标准误采用 Newey--West 异方差自相关稳健估计，
滞后阶数取 4。

\subsection{描述性统计}

表~\ref{tab:desc} 列出主要变量的分布。政策指引段情绪（tone\_P）与宏观回顾段情绪
（tone\_R）的相关系数为 0.375，二者相关但并非同一信息的重复度量，这为后续的
分段检验提供了必要的前提。事件日市场收益的均值为 19 个基点，标准差 186 个基点；
日内振幅均值为 2.09\%，意味着报告发布当日最高价与最低价之差平均约为前收盘价的
2\%。事前波动（发布前 60 个交易日市场收益的标准差）均值为 1.46\%。

\begin{table}[htbp]
\centering
\small
\begin{threeparttable}
\caption{描述性统计}
\label{tab:desc}
\begin{tabular}{lcccccc}
\toprule
变量 & 均值 & 标准差 & 最小值 & 中位数 & 最大值 & 观测数 \\
\midrule
政策指引段情绪（tone\_P，标准化） & 0.000 & 1.005 & $-2.773$ & 0.029 & 2.558 & 99 \\
宏观回顾段情绪（tone\_R，标准化） & 0.000 & 1.005 & $-2.928$ & 0.068 & 2.535 & 102 \\
发布后一日收益（\%）              & 0.19  & 1.86  & $-7.60$ & 0.30  & 4.63  & 102 \\
发布后五日收益（\%）              & 0.52  & 3.59  & $-12.59$ & 0.71 & 11.58 & 102 \\
事件日日内振幅（\%）              & 2.09  & 0.92  & 0.80  & 1.88  & 4.76  & 101 \\
事件日绝对收益（\%）              & 0.97  & 0.92  & 0.00  & 0.66  & 3.70  & 102 \\
事前 60 日波动（\%）              & 1.46  & 0.64  & 0.61  & 1.33  & 3.51  & 102 \\
\bottomrule
\end{tabular}
\begin{tablenotes}[flushleft]
\footnotesize
\item 样本为 2001Q1--2026Q2 共 102 期报告；收益与振幅由 16 个申万一级行业的等权组合计算。
\end{tablenotes}
\end{threeparttable}
\end{table}

值得先说明的是情绪与市场波动状态的原始相关性。tone\_P 与事件日日内振幅的相关系数为
$-0.302$，与事件日绝对收益为 $-0.254$，与事前波动的相关系数为 $-0.232$。最后一个
数字提示，情绪较高的报告往往出现在市场相对平静的时期，若不加以控制，很容易把
``情绪降低波动''与``情绪在高波动期偏低''混为一谈。因此在回归中同时报告含事前波动
控制的设定，并以此作为主要判断依据。

\subsection{政策指引情绪与市场波动}

表~\ref{tab:vol} 是本章的主结果。第 (1) 列以事件日绝对收益衡量波动，tone\_P 的系数为
$-23.4$ 个基点（$t=-2.53$）；加入事前波动后（第 (2) 列），系数降至 $-16.6$，显著性
减弱至 10\% 水平，说明有一部分关联确实来自波动状态的持续性。由于绝对收益无法区分
``价格单边移动''与``日内大幅震荡''，第 (3)、(4) 列改用日内振幅，系数分别为 $-27.8$
（$t=-2.58$）与 $-18.5$（$t=-2.43$），控制事前波动后依然在 5\% 水平显著。

第 (5) 列同时纳入回顾段情绪。tone\_R 的系数为 $-2.1$（$t=-0.20$），几乎为零，而
tone\_P 维持在 $-17.8$（$t=-2.27$）。这一对比有明确含义：显著影响波动的并不是
``整篇报告听起来是否乐观''，而是其中面向未来的政策指引部分。第 (6) 列把振幅扩展到
发布后五个交易日，系数为 $-12.1$（$t=-2.03$），效应并未随窗口延长而累积，更接近于
一次性的信息释放而非持续影响。

就经济意义而言，以振幅均值 2.09\% 为参照，$-18.5$ 个基点相当于振幅下降约 9\%；
以标准差衡量，一个标准差的情绪变动对应约 0.2 个标准差的振幅变化。这一量级不算大，
但在 99 个事件、日频噪声显著的条件下已属不易识别。

需要提醒的是，这一幅度对是否包含危机年份较为敏感。剔除 2007--2010 年后，系数降至
$-15.5$（$t=-2.07$），约为基准的一半（见 5.5 节）。方向与显著性不变，但把 $-28$
读作一个稳定不变的效应量并不恰当，更稳妥的表述是：效应量在 15 至 28 个基点之间，
取决于样本是否包含金融危机时期。

分组视角补充了一个细节。把情绪按三分位分组后，低情绪组的平均振幅为 2.53\%，中、高
两组分别为 1.75\% 与 1.99\%（图 5.2）。若把样本切成五组，均值依次为
2.83\%、2.07\%、1.65\%、2.13\%、1.78\%，同样不是单调下降。也就是说，效应主要来自
情绪偏低的一端：措辞异常谨慎的时期波动显著更高，而中高情绪之间并无梯度。这与
``沟通收敛分歧''的解释相容——市场在意的未必是措辞有多乐观，而是它是否透露出不安。

\begin{table}[htbp]
\centering
\small
\begin{threeparttable}
\caption{政策指引段情绪与市场波动}
\label{tab:vol}
\begin{tabular}{lcccccc}
\toprule
 & (1) & (2) & (3) & (4) & (5) & (6) \\
 & $|r|$ & $|r|$ & 振幅 & 振幅 & 振幅 & 五日振幅 \\
\midrule
tone\_P        & $-23.4^{**}$ & $-16.6^{*}$ & $-27.8^{***}$ & $-18.5^{**}$ & $-17.8^{**}$ & $-12.1^{**}$ \\
                & $(-2.53)$ & $(-1.96)$ & $(-2.58)$ & $(-2.43)$ & $(-2.27)$ & $(-2.03)$ \\
tone\_R        & & & & & $-2.1$ & \\
                & & & & & $(-0.20)$ & \\
事前波动        & & 46.0$^{***}$ & & 61.9$^{***}$ & 61.1$^{***}$ & 80.0$^{***}$ \\
                & & $(4.10)$ & & $(4.24)$ & $(4.14)$ & $(4.16)$ \\
\midrule
观测数          & 99 & 99 & 98 & 98 & 98 & 99 \\
$R^2$           & 0.065 & 0.162 & 0.091 & 0.267 & 0.268 & 0.355 \\
\bottomrule
\end{tabular}
\begin{tablenotes}[flushleft]
\footnotesize
\item 被解释变量为事件日市场波动，单位为百分比；情绪与事前波动的系数均已换算为基点
（1 个基点 $=0.01$ 个百分点），其中事前波动的系数读作``事前波动每上升 1 个百分点''。
括号内为 Newey--West 稳健 $t$ 值。$^{*}$、$^{**}$、$^{***}$ 分别表示 10\%、5\%、1\% 显著性水平。
\end{tablenotes}
\end{threeparttable}
\end{table}

\subsection{制度断点：2014 年前后}

样本期内货币政策框架发生了明显变化：2013 年流动性紧张、2014 年利率市场化加速、
2015 年汇率形成机制调整，报告自身的章节结构也在 2015 年前后改写。这些变化提示
情绪的作用方式未必稳定。表~\ref{tab:break} 在方程中引入 2014 年虚拟变量及其与
情绪的交互项，直接检验这一可能。

收益方程中，交互项系数为 $-75.6$ 个基点（$t=-2.11$），五日窗口扩大到 $-181.4$
（$t=-2.65$），即 2014 年以前情绪对收益方向无显著影响（点估计 $+31.6$，$t=0.97$），
2014 年以后转为负向。波动方程则呈现相反的变化：2014 年以前的系数为 $-35.0$
（$t=-3.10$），交互项为 $+26.6$（$t=1.59$），说明波动渠道在后期减弱，但未完全消失。

需要说明的是断点的选择并非由数据挑选而来，而是对应上述制度节点。为检验这一处理是否
稳健，表~\ref{tab:breaksens} 将断点依次替换为 2008、2010、2012、2014、2016、2018 年。
2008 年作断点时交互项接近于零，其余各年份的系数方向一致，其中 2010、2012、2014、2018
年显著，2016 年偏弱。这说明变化是渐进的，把它归因于某一个特定年份并不合适，但
``2010 年代之后效应改变''这一判断是站得住的。

\begin{table}[htbp]
\centering
\small
\begin{threeparttable}
\caption{制度断点检验}
\label{tab:break}
\begin{tabular}{lcccc}
\toprule
 & (1) & (2) & (3) & (4) \\
 & 一日收益 & 五日收益 & 日内振幅 & $|r|$ \\
\midrule
tone\_P              & $+31.6$   & $+73.7$   & $-35.0^{***}$ & $-32.7^{***}$ \\
                      & $(0.97)$  & $(1.22)$  & $(-3.10)$     & $(-2.97)$ \\
tone\_P $\times$ 2014后 & $-75.6^{**}$ & $-181.4^{***}$ & $+26.6$ & $+25.1$ \\
                      & $(-2.11)$ & $(-2.65)$ & $(1.59)$      & $(1.53)$ \\
\midrule
观测数                & 99 & 99 & 98 & 99 \\
$R^2$                 & 0.076 & 0.077 & 0.190 & 0.084 \\
\bottomrule
\end{tabular}
\begin{tablenotes}[flushleft]
\footnotesize
\item 各列均控制 2014 年虚拟变量本身；系数单位为基点，括号内为 Newey--West 稳健 $t$ 值。
\end{tablenotes}
\end{threeparttable}
\end{table}

\begin{table}[htbp]
\centering
\small
\begin{threeparttable}
\caption{断点年份敏感性（交互项系数与 $t$ 值）}
\label{tab:breaksens}
\begin{tabular}{lccc}
\toprule
断点年份 & 一日收益 & 五日收益 & 日内振幅 \\
\midrule
2008 & $-5.9\ (-0.14)$   & $-1.6\ (-0.02)$    & $+1.8\ (0.10)$ \\
2010 & $-86.6^{**}\ (-2.41)$ & $-193.1^{***}\ (-2.69)$ & $+34.5^{**}\ (2.04)$ \\
2012 & $-76.6^{**}\ (-2.20)$ & $-183.5^{***}\ (-2.74)$ & $+27.4^{*}\ (1.71)$ \\
2014 & $-75.6^{**}\ (-2.11)$ & $-181.4^{***}\ (-2.65)$ & $+26.6\ (1.59)$ \\
2016 & $-57.7\ (-1.47)$  & $-134.9^{*}\ (-1.88)$   & $+34.4^{**}\ (2.07)$ \\
2018 & $-64.4^{*}\ (-1.82)$ & $-164.9^{**}\ (-2.29)$ & $+16.2\ (0.79)$ \\
\bottomrule
\end{tabular}
\begin{tablenotes}[flushleft]
\footnotesize
\item 每行对应一次单独的回归，仅改变断点年份，其余设定与表~\ref{tab:break} 相同。
\end{tablenotes}
\end{threeparttable}
\end{table}

\subsection{2014 年之后的方向性效应}

既然断点检验显示后期情绪与收益转为负相关，这里单独报告 2014 年以后的子样本。
表~\ref{tab:post14} 显示，tone\_P 对发布后一日收益的系数为 $-44.0$ 个基点
（$t=-2.89$），五日窗口为 $-107.8$（$t=-3.30$）；同时，波动类方程在子样本中
全部不显著。两组结果互为印证：后期的定价体现在方向上，而不在波动上。

对这一结果需要保持克制。首先，负号的含义并非``情绪越宽松市场越差''这么简单，
更可能反映宽松措辞集中出现在基本面走弱的时点，市场将政策宽松解读为经济前景的
负面信号。其次，该结果对危机年份高度敏感：剔除 2015 年后，一日窗口的交互项由
$-75.6$（$t=-2.11$）降至 $-59.0$（$t=-1.62$）；剔除 2008--2009 年后降至 $-45.5$
（$t=-1.56$）。任一段被剔除，5\% 水平的显著性即不复存在。因此本文不把它作为独立
结论，仅记录符号的方向性变化。

\begin{table}[htbp]
\centering
\small
\begin{threeparttable}
\caption{2014 年以后的子样本}
\label{tab:post14}
\begin{tabular}{lcccc}
\toprule
 & (1) & (2) & (3) & (4) \\
 & 一日收益 & 五日收益 & 一日·双段 & 日内振幅 \\
\midrule
tone\_P & $-44.0^{***}$ & $-107.8^{***}$ & $-47.2^{**}$ & $-8.4$ \\
        & $(-2.89)$     & $(-3.30)$      & $(-2.36)$    & $(-0.68)$ \\
tone\_R &               &                & $+5.9$       & \\
        &               &                & $(0.28)$     & \\
\midrule
观测数  & 50 & 50 & 50 & 49 \\
\bottomrule
\end{tabular}
\begin{tablenotes}[flushleft]
\footnotesize
\item 子样本为 2014Q1--2026Q2。系数单位为基点，括号内为 Newey--West 稳健 $t$ 值。
\end{tablenotes}
\end{threeparttable}
\end{table}

\subsection{稳健性检验}

表~\ref{tab:robust} 汇总四个方面的检验，表~\ref{tab:robust2} 补充三项针对性检验。

\emph{市场口径}。把等权组合由 16 个行业扩展为 28 个行业，主结果基本不变：
日内振幅系数 $-26.7$（$t=-2.58$），断点交互项 $-73.1$（$t=-2.01$），五日交互项
$-183.6$（$t=-2.71$）。结论不依赖于行业集合的选择。

\emph{词典}。以本文自建的货币政策领域词典（积极 159 词、消极 139 词）重新度量情绪，
各系数方向与主结果一致，但量级缩小、显著性下降：绝对收益方程仍显著（$-19.1$，
$t=-2.43$），日内振幅方程降至 $-7.8$（$t=-0.66$），断点交互项降至 $-60.5$
（$t=-1.25$）。该词典每千字仅命中约 68 个情感词，主词典为 104 个，覆盖率差异恰好
对应系数的等比例收缩，符合测量误差导致衰减的预期。因此这一结果是``方向稳健、精度不足''，
而非``结论被推翻''。

\emph{影响点}。逐一剔除 2015 年（股市异常波动）与 2020 年（疫情冲击），并对收益做
1\%/99\% 缩尾。波动类结果对上述处理均不敏感，日内振幅系数稳定在 $-25.8$ 至 $-27.9$
之间。断点交互项则较为敏感：剔除 2015 年后，一日窗口的交互项由 $-75.6$（$t=-2.11$）
降至 $-59.0$（$t=-1.62$），显著性消失；五日窗口降至 $-152.5$（$t=-2.29$），仍在 5\%
水平显著。剔除 2020 年与收益缩尾后，一日交互项分别为 $-73.0$ 与 $-74.4$，保持在 5\%
水平。可见后期方向效应确有相当一部分由 2015 年的数期报告贡献，这解释了上文为何
采取保留态度，也说明该效应在更长窗口上比在单日窗口上更稳定。

\emph{置换检验}。将情绪序列与事件随机配对 500 次，观察真实 $t$ 值在随机分布中的位置。
日内振幅方程的置换 $p$ 值为 0.020，加入回顾段与事前波动后为 0.068，断点交互项为 0.058。
置换检验不依赖分布假设，其结果与稳健 $t$ 检验基本一致：波动结果在 5\% 水平成立，
断点结果处于边际。

\emph{危机期敏感性}。2007--2009 年的措辞与波动同时处于极端，是本文最担心的混淆来源。
剔除 2007--2010 年后，振幅方程系数降至 $-15.5$（$t=-2.07$），剔除 2008--2009 年后为
$-18.8$（$t=-2.34$），加入年份趋势后为 $-20.7$（$t=-2.18$）。方向与显著性保持不变，
幅度约为基准的一半。断点交互项则不同：剔除 2008--2009 年后降至 $-45.5$（$t=-1.56$），
不再显著。可见波动结果稳健但幅度受危机年份影响，方向结果则更依赖少数极端时期。

\emph{安慰剂窗口}。把事件日整体平移 $\pm$5、$\pm$10、$\pm$20 个交易日重估。平移 $\pm$10
天以上时，振幅方程的系数不再显著（$t$ 值均小于 1.1），断点交互项也不再成立，说明效应
与发布时点相关。但平移 $-5$ 天时系数仍有 $-26.6$（$t=-1.79$），保留了约九成的幅度。
这表明效应并非严格的单日冲击，而是集中在发布前后约一周的窗口内，与报告部分内容提前
释放的解释相容。

\emph{政策操作控制}。一个自然的担心是，文本情绪只是货币政策操作的代理。本文整理了三类
总量操作——全面降准 24 次、1 年期 LPR 下调 12 次、7 天逆回购利率下调 12 次——按报告期
汇总为季度操作强度，加入振幅方程。情绪系数由 $-27.8$ 变为 $-27.3$（$t=-2.53$），加入
事前波动后为 $-16.7$（$t=-2.17$）；三类操作本身均不显著，与情绪的相关性也只在 $-0.01$
到 $-0.21$ 之间。断点交互项由 $-75.5$ 变为 $-73.9$（$t=-2.12$）。这说明波动效应并非
由降准降息驱动。

\begin{table}[htbp]
\centering
\small
\begin{threeparttable}
\caption{稳健性检验汇总}
\label{tab:robust}
\begin{tabular}{lcccc}
\toprule
 & 主结果 & 28 行业 & 替代词典 & 剔除 2015 \\
\midrule
\multicolumn{5}{l}{\emph{面板 A：日内振幅（系数，$t$ 值）}} \\
tone\_P & $-27.8^{***}$ & $-26.7^{***}$ & $-7.8$ & $-25.8^{**}$ \\
        & $(-2.58)$ & $(-2.58)$ & $(-0.66)$ & $(-2.19)$ \\
\midrule
\multicolumn{5}{l}{\emph{面板 B：断点交互项（系数，$t$ 值）}} \\
tone\_P $\times$ 2014后（一日） & $-75.6^{**}$ & $-73.1^{**}$ & $-60.5$ & $-59.0$ \\
        & $(-2.11)$ & $(-2.01)$ & $(-1.25)$ & $(-1.62)$ \\
tone\_P $\times$ 2014后（五日） & $-181.4^{***}$ & $-183.6^{***}$ & $-109.3$ & $-152.5^{**}$ \\
        & $(-2.65)$ & $(-2.71)$ & $(-1.25)$ & $(-2.29)$ \\
\midrule
\multicolumn{5}{l}{\emph{面板 C：置换检验 $p$ 值}} \\
 & 0.020 & — & — & — \\
\bottomrule
\end{tabular}
\begin{tablenotes}[flushleft]
\footnotesize
\item ``替代词典''指以自建货币政策领域词典重新计算情绪指标的结果；``剔除 2015''指去掉 2015 年
四期报告后重估。面板 C 仅报告日内振幅方程的置换 $p$ 值。
\end{tablenotes}
\end{threeparttable}
\end{table}

\begin{table}[htbp]
\centering
\small
\begin{threeparttable}
\caption{补充检验：危机期、安慰剂窗口与政策操作}
\label{tab:robust2}
\begin{tabular}{lcc}
\toprule
设定 & 日内振幅（系数，$t$ 值） & 断点交互项（系数，$t$ 值） \\
\midrule
基准 & $-27.8^{***}$ $(-2.58)$ & $-75.5^{**}$ $(-2.11)$ \\
剔除 2007--2010 & $-15.5^{**}$ $(-2.07)$ & $-36.6$ $(-1.27)$ \\
剔除 2008--2009 & $-18.8^{**}$ $(-2.34)$ & $-45.5$ $(-1.56)$ \\
加入年份趋势 & $-20.7^{**}$ $(-2.18)$ & — \\
事件日前移 20 个交易日 & $+8.5$ $(0.46)$ & $+9.6$ $(0.40)$ \\
事件日后移 20 个交易日 & $-12.5$ $(-1.21)$ & $-35.1$ $(-1.51)$ \\
控制三类政策操作 & $-27.3^{**}$ $(-2.53)$ & $-73.9^{**}$ $(-2.12)$ \\
\bottomrule
\end{tabular}
\begin{tablenotes}[flushleft]
\footnotesize
\item 系数单位为基点，括号内为 Newey--West 稳健 $t$ 值。安慰剂检验按平移后的日期
重新计算振幅与收益，情绪变量不变。政策操作指全面降准、1 年期 LPR 与 7 天逆回购利率，
按报告期汇总为季度强度。``—''表示该项未单独估计。
\end{tablenotes}
\end{threeparttable}
\end{table}

综合来看，本章的证据支持一个有限但明确的判断：政策指引段的文本情绪确实影响市场，
但其作用首先体现在波动而非收益方向上；方向性效应在 2014 年之后才出现，且稳健性
弱于波动结果。这一区分对理解央行沟通的作用机制是重要的——若只报告全样本的平均效应，
两个方向相反的结果会相互抵消，从而得出``沟通无效''的错误结论。
